\section{Bounded paths exclude AUC renormings}\label{sec:obstruction}
The midpoint arguments depend on the multiplier distribution. Here we
use only positivity, a nonconstant mean-one law, and a finite second moment
to rule out AUC renormings. At every vertex the child increments are
weakly null and have the same positive norm. An AUC norm would therefore
grow by a fixed factor at each step of a suitably chosen branch, whereas
all path products have inherited norm one.

For each fixed vertex $v$ and its children $v^\frown n$, put
\[
 D_{v,n}=P_{v^\frown n}-P_v=P_v(W_{v^\frown n}-1).
\]
Independence and~\eqref{prob:product-moments} give
\begin{equation}\label{prob:child-moments}
 \norm{D_{v,n}}=\rho:=\E|W-1|>0,\qquad
 \E D_{v,n}^2=m_2^{|v|}(m_2-1).
\end{equation}
For fixed $v$ the sequence $(D_{v,n})_n$ is pairwise orthogonal in $L^2$
and has constant finite $L^2$ norm.  Bessel's inequality therefore makes
it weakly null in $L^2$.  The inclusion $L^2\longrightarrow L^1$ is
continuous on this probability space.  Moreover, each continuous linear
functional on $X_W$ extends to $L^1$ by Hahn--Banach and hence restricts
to a continuous functional on $L^2$.  Thus $(D_{v,n})_n$ is weakly null
in $X_W$ itself.

Suppose $N$ is equivalent to the inherited norm, with
$a\norm{z}\le N(z)\le b\norm{z}$ for $0<a\le b<\infty$.
We will show the precise obstruction
\begin{equation}\label{eq:auc-zero}
 \overline\delta_N\left(\frac{a\rho}{2b}\right)=0.
\end{equation}
The usual one-sided modulus is nonnegative: for each unit center $x$, a
norming functional gives a finite-codimensional kernel on which
$N(x+ty)\ge1$. It remains to rule out positivity at the displayed scale.

Assume, to the contrary, that
\[
 t_0=\frac{a\rho}{2b},\qquad
 \delta=\overline\delta_N(t_0)>0,\qquad \gamma=\delta/2.
\]
At a vertex $v$, we have $a\le N(P_v)\le b$. The modulus at the unit center
$P_v/N(P_v)$ supplies a closed finite-codimensional subspace $F$ such that
\[
 N(P_v+z)\ge(1+\gamma)N(P_v)
 \quad(z\in F,\ N(z)\ge t_0N(P_v)).
\]
Indeed, its definition first gives the estimate at equality of the radius.
For $N(x)=N(w)=1$, with $w\in F$ and $r\ge t_0$, convexity yields
\[
 N(x+rw)-1\ge\frac r{t_0}\bigl(N(x+t_0w)-1\bigr),
\]
so the estimate extends to all larger radii.

Equivalent norms induce the same weak topology. The quotient images of the
weakly null sequence $D_{v,n}$ in the finite-dimensional space $X_W/F$
therefore tend to zero in quotient norm. We may choose $z_n\in F$ with
$N(D_{v,n}-z_n)\to0$. Eventually
\[
 N(z_n)\ge a\rho/2=t_0b\ge t_0N(P_v).
\]
Choose such an $n$ with approximation error at most $\gamma N(P_v)/2$.
Then
\[
 N(P_{v^\frown n})
 \ge N(P_v+z_n)-N(D_{v,n}-z_n)
 \ge(1+\gamma/2)N(P_v).
\]
The same $\gamma$ works at every vertex, although $F$ and $n$ may change.
Repeatedly selecting a child gives a depth-$k$ product with norm at least
$a(1+\delta/4)^k$, contrary to the uniform upper bound $b$.
This proves~\eqref{eq:auc-zero} and the renorming assertion of
Theorem~\ref{thm:main}.

Finally $0<\rho<2$: nonconstancy gives positivity, and
$\E|W-1|=2-2\E\min\{W,1\}<2$ because $W>0$ almost surely.
Thus $0<t_0<1$, so positivity of the one-sided modulus on $(0,1)$ alone
would already contradict the tree. For exponential multipliers
$\rho=2/e$; indeed $\E|Z-1|=2\E(Z-1)_+=2\int_1^\infty e^{-s}\,ds$.
The weak-null calculation requires only a finite $L^2$ bound at each
fixed parent, not a bound uniform over the depth.
